% Convergence of three collocation schemes on test problems P1--P6.
% Compile this file on its own (pdflatex collocationConvergence.tex) to produce
% collocationConvergence.pdf, which the paper includes. It is set in the same
% fonts as the paper (Baskervaldx with newtx maths), at the paper's text width,
% so that it is included at its natural size; it has no background.
% The data in data/*.csv (columns: n, error) were produced by the MATLAB
% scripts of the paper.
\documentclass{standalone}
\usepackage[lining]{Baskervaldx}
\usepackage[baskervaldx,vvarbb,bigdelims]{newtxmath}
\usepackage{pgfplots}
\usepgfplotslibrary{groupplots}
\pgfplotsset{compat=1.18}

% MATLAB's default line colours, as in the original figure
\definecolor{c1}{RGB}{0,114,189}
\definecolor{c2}{RGB}{217,83,25}
\definecolor{c3}{RGB}{237,177,32}
\definecolor{c4}{RGB}{126,47,142}
\definecolor{c5}{RGB}{119,172,48}
\definecolor{c6}{RGB}{77,190,238}

% Text width of the paper (siamonline and daarticle)
\newlength{\figwidth}\setlength{\figwidth}{444.53284pt}

\pgfplotsset{
  conv/.style={
    xmode=log, ymode=log,
    ymin=1e-7, ymax=1e3, ytick={1e-5,1e0},
    height=4.6cm, scale only axis,
    grid=both, major grid style={black!10}, minor grid style={black!5},
    tick style={black!70}, tick label style={font=\footnotesize},
    title style={font=\footnotesize, align=center},
    every axis plot/.append style={line width=.6pt, mark=*, mark size=.7pt},
    cycle list={{c1},{c2},{c3},{c4},{c5},{c6}},
    xlabel={$n$}, xlabel style={font=\footnotesize, yshift=3pt},
  },
}

\newcommand{\convplots}[1]{%
  \foreach \p in {P1,P2,P3,P4,P5,P6}{%
    \addplot table[col sep=comma, x index=0, y index=1] {data/#1-\p.csv};}}

\begin{document}
\begin{tikzpicture}
  % Panel widths are proportional to the number of decades shown, so that a
  % decade has the same length in every panel
  \pgfmathsetlengthmacro{\decade}{(\figwidth-70pt)/(3+2.477+2)}
  \begin{groupplot}[
      group style={group size=3 by 1, horizontal sep=10pt, yticklabels at=edge left},
      conv]
    \nextgroupplot[
      width=3*\decade, xmin=1, xmax=1e3,
      % no label at 10^3, which would collide with the next panel's 10^0
      xtick={1,10,100,1000}, xticklabels={$10^0$,$10^1$,$10^2$,},
      title={Finite Element Collocation\\(Trapezium)},
      ylabel={$\displaystyle\max_k \| u(t_k) - u_n(t_k) \|_\infty$},
      ylabel style={font=\footnotesize},
      legend columns=1, legend cell align=left,
      legend style={font=\footnotesize, draw=black!20, fill=none,
        at={(.98,.98)}, anchor=north east, row sep=-2pt},
      ]
      \convplots{FEColl}
      \legend{P1,P2,P3,P4,P5,P6}
    \nextgroupplot[
      width=2.477*\decade, xmin=1, xmax=300,
      title={Spectral Collocation\\(Trapezium)}]
      \convplots{SpecCollTrap}
    \nextgroupplot[
      width=2*\decade, xmin=1, xmax=1e2,
      title={Spectral Collocation\\(Clenshaw--Curtis)}]
      \convplots{SpecCollCC}
  \end{groupplot}
  % Panel labels
  \foreach \k/\l in {1/a,2/b,3/c}{%
    \node[font=\footnotesize, anchor=north] at ([yshift=-17pt]group c\k r1.south) {(\l)};}
\end{tikzpicture}
\end{document}
